\section{Model C: cofactor depletion and bistability}
%==========================================================================

\subsection{What is added}

Models A and B treat alternative clearance as a fixed capacity. It is not: elimination
proceeds by conjugating propionate with carnitine, and \textbf{each molecule eliminated
consumes one molecule of carnitine}, which is excreted with the product. Carnitine is a
consumed reagent, not a catalyst. Secondary carnitine deficiency is a documented feature of
the disease and the reason carnitine supplementation is standard therapy.

Promoting carnitine to a state variable:

\begin{equation}
\begin{aligned}
\dot m &= s(t) \;-\; \theta\, g(m) \;-\; k\, c\, m, \\
\dot c &= R \;-\; k\, c\, m,
\end{aligned}
\label{eq:C}
\end{equation}

with $c(t)$ available carnitine and $R$ its replenishment rate. The term $k c m$ appears in
both equations with the same coefficient, expressing the one-to-one stoichiometry. The
enzymatic term $\theta g(m)$ is unchanged and remains Michaelis-Menten, with $\theta$ the
residual $V_{\max}$ and $K_m$ the Michaelis constant; it is only the alternative channel
that is re-specified.

\subsection{Consequence: positive feedback}

The loop is
\[
m \uparrow \;\Longrightarrow\; \text{carnitine consumed} \;\Longrightarrow\;
c \downarrow \;\Longrightarrow\; \text{clearance} \downarrow \;\Longrightarrow\; m \uparrow .
\]

If depletion outpaces replenishment, effective clearance \emph{decreases} with $m$ rather
than merely saturating. The clearance curve becomes non-monotone, and the system admits
\textbf{two} equilibria at moderate load: a stable one at low $m$, and an unstable one at
higher $m$. As load rises the two approach, collide and annihilate, which is a genuine
saddle-node.

\subsection{Why this would matter clinically}

The unstable equilibrium is a \emph{point of no return}. Beyond it, $\dot m > 0$ even at
baseline load, so the metabolite continues to rise after the precipitating illness has
resolved. Recovery requires active intervention to return the patient to the left of the
separatrix, which is precisely what intravenous glucose, carnitine and haemofiltration
accomplish.

This predicts the observed clinical fact that decompensations do not resolve spontaneously,
which Models A and B do not explain. An analogous depletion argument applies to free
coenzyme A.

\subsection{Why it is not the working model}

\begin{enumerate}[leftmargin=*, itemsep=2pt]
\item \textbf{Unverified.} Whether depletion is quantitatively strong enough to render
      clearance decreasing, rather than merely saturating, has not been established. The
      test is direct and requires no new experiment: free carnitine and the
      acylcarnitine/free-carnitine ratio are measured routinely, so one may ask whether free
      carnitine collapses during decompensation relative to the same patient's baseline.
\item \textbf{Identifiability.} $k_e$ is replaced by the product $k c(t)$ with $c$ a further
      unmeasured state and $R$ a further unknown. This worsens an already unfavourable
      parameter-to-measurement ratio.
\item \textbf{Incompleteness.} Equations \eqref{eq:C} omit free-carnitine excretion, renal
      reabsorption, and the dependence of propionylcarnitine elimination on renal function.
      Adding one state to capture half a mechanism may be worse than not adding it, since it
      confers the appearance of having modelled depletion without having done so.
\end{enumerate}

\subsection{Testable discriminant}

Models B and C differ in a clinically checkable prediction:

\begin{center}
\begin{tabular}{@{}ll@{}}
\toprule
\textbf{B} & the crisis resolves when the precipitating illness resolves \\
\textbf{C} & beyond the separatrix the crisis is self-sustaining; intervention is required \\
\bottomrule
\end{tabular}
\end{center}

Clinical practice is consistent with C, but this is an observation and not a demonstration.

\vspace{0.5em}
\hrule
\vspace{1em}

%==========================================================================
